\documentclass[12pt]{article}
\usepackage[utf8]{inputenc}
\usepackage[T1]{fontenc}
\usepackage[margin=1in]{geometry}
\usepackage{setspace}
\usepackage{xcolor}
\usepackage[colorlinks=true, urlcolor=blue!60!black, linkcolor=blue!60!black]{hyperref}
\setlength{\parindent}{2em}
\setlength{\parskip}{1.25em}
\renewcommand{\baselinestretch}{1.0}
\begin{document}
\begin{center}
{\LARGE \textit{Biden v.\ Nebraska}, 600 U.S.\ 477 (2023)\par}
\vspace{0.8em}
{\large Case Brief --- Four-Page Edition\par}
\end{center}
\vspace{1em}

\section{Memory Jogger}

The Court held 6--3 that Missouri had Article III standing through financial harm to its state-created loan servicer MOHELA, and that the Secretary of Education's HEROES Act authority to ``waive or modify'' student-loan provisions did not authorize a \$430 billion mass debt-cancellation program, under both ordinary statutory text and the Major Questions Doctrine (\href{https://www.law.cornell.edu/supremecourt/text/22-506}{600 U.S. 477}).

\section{Detailed Case Facts}

The Higher Education Act lets the Secretary discharge loan principal only in narrow circumstances: public service, death or disability, bankruptcy, and school closure or false certification. The HEROES Act of 2003 permits the Secretary to ``waive or modify any statutory or regulatory provision applicable to the student financial assistance programs under title IV'' in connection with a war, military operation, or national emergency, so that affected borrowers are not left financially worse off. Before 2020, waivers were modest and procedural.

After the COVID--19 emergency declaration, both Administrations paused repayment and interest. In October 2022, Secretary Cardona announced a plan discharging up to \$10,000 (or \$20,000 for Pell Grant recipients) for borrowers earning under \$125,000. The CBO estimated the cost at roughly \$430 billion, covering about 98.5\% of borrowers. Six States sued. Missouri relied on MOHELA, a public instrumentality created by state statute that serviced federal loans and would lose an estimated \$44 million a year in fees, even though MOHELA itself declined to join the suit.

\section{Procedural History}

The Eastern District of Missouri dismissed for lack of standing. The Eighth Circuit granted a nationwide injunction pending appeal. The Supreme Court granted certiorari before judgment and, on June 30, 2023, reversed and remanded; the district court later vacated the program. In the companion case, \textit{Department of Education v. Brown}, the Court unanimously held that two individual borrowers lacked standing.

\section{Judicial Votes}

Reversed and remanded, \textbf{6--3}.
\begin{itemize}
\item \textbf{Majority:} Chief Justice Roberts, joined by Justices Thomas, Alito, Gorsuch, Kavanaugh, and Barrett.
\item \textbf{Concurring:} Justice Barrett, joining in full and defending the Major Questions Doctrine as a contextual tool consistent with textualism.
\item \textbf{Dissenting:} Justice Kagan, joined by Justices Sotomayor and Jackson.
\end{itemize}

\section{Holding}

(1) When a State creates, governs, supervises, and may dissolve a public corporation performing an essential public function, a direct financial injury to that instrumentality is an injury to the State itself, even if the corporation can sue in its own name. (2) The power to ``waive or modify'' permits nullifying specific requirements or making modest, incremental adjustments, not creating a novel \$430 billion forgiveness regime. (3) Major Questions principles apply to benefits programs as well as regulatory mandates, so an agency claiming unheralded power of vast significance must point to clear congressional authorization. Footnote 5 expressly left open the legality of the COVID--19 payment pause.

\section{Analysis of Opinions}

\textbf{Majority (Roberts).} \emph{Standing:} the \$44 million loss to MOHELA was conceded. Missouri created MOHELA by statute, its board consists of state officials and gubernatorial appointees, it reports to the State, and only Missouri can dissolve it. Under \textit{Arkansas v. Texas} and \textit{Lebron v. National Railroad Passenger Corp.}, a State's use of a separate corporation does not bar it from suing for harm to that instrumentality. \emph{Text:} drawing on \textit{MCI v. AT\&T}, ``modify'' connotes moderate change and cannot authorize fundamental changes to Congress's scheme; the plan modified the Education Act only in the sense that the French Revolution modified the status of the French nobility. Nor was anything ``waived''; the Secretary substituted an elaborate new regime. \emph{Major questions:} under \textit{West Virginia v. EPA}, the plan's unheralded scale, staggering cost, political salience amid more than 80 failed forgiveness bills, and implication of the power of the purse demanded clear authorization, and there was none.

\textbf{Barrett (concurring).} Strong-form substantive canons sit in tension with textualism, so Barrett recast the doctrine as a linguistic tool. Her babysitter analogy: told to make sure the kids have fun, a sitter should not take them on a two-day amusement-park trip without explicit permission. Because Article I vests all legislative powers in Congress, constitutional structure is part of the context for reading delegations.

\textbf{Kagan (dissenting).} The States were ideological plaintiffs. Missouri's treasury would lose nothing because MOHELA is legally and financially independent, and MOHELA declined to sue; the majority turned standing into ``a lawyer's game.'' On the merits, ``waive or modify'' together authorize changes along a spectrum, and under her dirty-bomb hypothetical the Act plainly lets the Secretary cancel loans for displaced borrowers. Reliance on the major-questions doctrine was a tell that the text could not sustain the result.

\section{Examples: Future Applications}

\begin{enumerate}
\item \textbf{Same-Side (Invalid): Farm-Debt Cancellation.} An agriculture secretary uses a ``waive or modify'' emergency clause to cancel \$350 billion in farm-loan principal. A state-chartered lending authority's lost fees give its State standing, and the program fails on text and major questions.
\item \textbf{Same-Side (Invalid): Mortgage Principal Reduction.} HUD forgives \$50,000 of principal on all federally backed mortgages at a \$500 billion cost. Even within HUD's domain, ``waive or modify'' cannot invent sweeping debt cancellation Congress has declined to enact.
\item \textbf{Opposite-Side (Valid): Hurricane Waivers.} After a hurricane emergency, the Secretary extends deadlines and relaxes service rules for borrowers in disaster counties, mirroring the modest 2003 and 2012 waivers the Court endorsed.
\item \textbf{Opposite-Side (Valid): Express Authorization.} Congress expressly authorizes emergency cancellation of up to \$15,000 per borrower under a \$200 billion cap; that is clear congressional authorization.
\item \textbf{Fence-Sitter: A Three-Year Payment Pause.} The Secretary cancels no principal but suspends payments and sets interest at zero for three years, at a \$150 billion cost. It looks like a temporary waiver and the CARES Act arguably ratified the pause, but its magnitude rivals what triggered the doctrine here, and Footnote 5 left the question open.
\end{enumerate}

\section{Critique}

\textbf{Progressive and institutionalist critiques.} Critics echoed Kagan that letting Missouri sue for an unwilling, financially insulated corporation eroded Article III's injury requirement. The Tenth Circuit later held in \textit{Good v. U.S. Department of Education} that MOHELA is not an arm of the State for sovereign immunity, producing an asymmetry. Jed Handelsman Shugerman argued that the majority used a purposivist Major Questions Doctrine to rewrite broad text. Kevin Tobia, Daniel Walters, and Brian Slocum tested Barrett's analogies empirically and found that ordinary people do not uniformly read open-ended delegations to exclude high-stakes actions.

\textbf{Conservative and originalist critiques.} William Baude and Samuel Bray argued that the MOHELA theory stretched traditional proper-party limits. Chad Squitieri and Adrian Vermeule warned that Barrett's contextual approach risks collapsing textualism into purposivism, while strict nondelegation originalists would enforce Article I directly rather than inventing clear-statement rules.

\section{Key Quotations}

\begin{quote}
``We hold today that the Act allows the Secretary to `waive or modify' existing statutory or regulatory provisions applicable to financial assistance programs under the Education Act, not to rewrite that statute from the ground up.''
\end{quote}
\noindent --- Chief Justice Roberts, Opinion of the Court, 600 U.S. 477.

\begin{quote}
``Rather, I understand it to emphasize the importance of context when a court interprets a delegation to an administrative agency. Seen in this light, the major questions doctrine is a tool for discerning---not departing from---the text's most natural interpretation.''
\end{quote}
\noindent --- Justice Barrett, concurring, 600 U.S. 477.

\begin{quote}
``The dissent is correct that this is a case about one branch of government arrogating to itself power belonging to another. But it is the Executive seizing the power of the Legislature.''
\end{quote}
\noindent --- Chief Justice Roberts, Opinion of the Court, 600 U.S. 477.

\end{document}
